\section{The critical background and finite-volume comparison}
\label{flow:section}

We now choose the temperature by the decay of the homogeneous background.
The interpretation of this branch as critical will follow from the
correlations; no property of an interacting critical point is assumed
in its construction.  All maps in this section use the critical
nearest-neighbor conditional kernels of \Cref{model:reference}.

\subsection{A decaying homogeneous trajectory}

Choose the contraction parameter in \Cref{rg:linear} and numbers
\begin{equation}\label{flow:rates}
                 0<\rho<\gamma_0<\gamma<1.
\end{equation}
Fix the corresponding sufficiently large block ratio $L$ and initial
radius $r_0$, with $r_0$ also larger than the range of the fixed potential.
The choices may depend on $J,V$, but not on a domain or on the number of
insertions.  We may increase $r_0$ to meet all preceding calibration and
buffer thresholds.

Let $\mathcal H_j$ be the space of homogeneous, square-symmetric,
spin-flip-even background lists at scale $r_j$.  Use the uniformly
bounded splitting of \Cref{rg:linear} to write a background as
\[
                 h=t\boldsymbol v_j+w,\qquad w\in W_j.
\]
The norm $|t|+\norm{w}_{r_j}$ is uniformly equivalent to
$\norm{h}_{r_j}$.  Here $t$ is the coefficient of the effective thermal
list; its relation to the microscopic inverse temperature will be
determined by the preliminary transformation below.  In these
coordinates the analytic local map is
\begin{equation}\label{flow:coordinate-map}
\begin{split}
 t_{j+1}&=Lt_j+B_jw_j+n_j^1(t_j,w_j),\\
 w_{j+1}&=D_jw_j+n_j^2(t_j,w_j),
\end{split}
\end{equation}
where $\norm{D_j}\le\rho$, $\norm{B_j}\le C_L$, and the remainders
are analytic on a common complex neighborhood of zero.  They vanish
to second order and satisfy uniform quadratic and local Lipschitz
bounds there, by \Cref{rg:map}.

\begin{lemma}[The decaying graph]\label{flow:stable-graph}
There is $\epsilon>0$ such that, for every sufficiently small
$w_0\in W_0$, there is exactly one trajectory of
\eqref{flow:coordinate-map} with this initial stable coordinate,
$\sup_j(|t_j|+\norm{w_j}_{r_j})<\epsilon$, and
\[
        \sup_{j\ge0}\gamma_0^{-j}
                (|t_j|+\norm{w_j}_{r_j})<\infty.
\]
It satisfies
\begin{equation}\label{flow:stable-bound}
 |t_j|+\norm{w_j}_{r_j}\le C\norm{w_0}_{r_0}\gamma_0^j.
\end{equation}
Its initial thermal coordinate $t_0=\psi(w_0)$ is analytic, with
$\psi(0)=0$.  This graph is transverse to the thermal line.
\end{lemma}

\begin{proof}
We give the forward-stable, backward-thermal solution operator, since
the spaces and linear maps depend on the scale. Sequence-space
constructions for scale-dependent renormalization dynamics are also
developed in \cite{BBS2015}; the contraction needed here is proved
explicitly below.  For an imposed forcing $f_j=(f_j^1,f_j^2)$, first solve
the stable equation forward:
\[
 w_j=D_{j-1}\cdots D_0w_0+
       \sum_{i=0}^{j-1}D_{j-1}\cdots D_{i+1}f_i^2,
\]
where an empty product is the identity.  Then solve the thermal
equation backward:
\begin{equation}\label{flow:backward}
 t_j=-\sum_{i=j}^{\infty}L^{j-1-i}(B_iw_i+f_i^1).
\end{equation}
On the sequence space with norm
\[
 \norm{(t,w)}_{\gamma_0}
   =\sup_{j\ge0}\gamma_0^{-j}(|t_j|+\norm{w_j}_{r_j}),
\]
these operators are bounded.  Indeed the stable forcing sum costs
at most $(\gamma_0-\rho)^{-1}\norm{f^2}_{\gamma_0}$, and the
backward sum costs at most
$(L-\gamma_0)^{-1}(C_L\norm{w}_{\gamma_0}
+\norm{f^1}_{\gamma_0})$.  Any solution in this space must obey
\eqref{flow:backward}: iterate its thermal recurrence to time $k$
and let $k\to\infty$, using $L^{j-k}t_k\to0$.

Substitute $f_j=n_j(t_j,w_j)$ in this bounded linear solution operator.
The uniform quadratic bounds imply
\[
 \norm{n(t,w)}_{\gamma_0}
       \le C_L\norm{(t,w)}_{\gamma_0}^2,
\]
and its Lipschitz constant on a sequence ball of radius $s$ is at most
$C_Ls$.  A sufficiently small ball, with radius a fixed multiple of
$\norm{w_0}_{r_0}$, is therefore invariant and contractive.  The
contraction theorem proves existence, uniqueness, and
\eqref{flow:stable-bound} in this weighted sequence ball.  To obtain
the stated uniqueness, any two trajectories in the ordinary
$\epsilon$-ball with finite weighted norm obey the same linear
solution formulas.  Their difference is bounded in weighted norm by
$C\epsilon$ times itself, by the local Lipschitz bound on the
remainders and the bounded solution operator.  Choose $\epsilon$ so
that $C\epsilon<1$; the constructed trajectory lies in this ball
when $w_0$ is sufficiently small.

The complex analytic contraction, or
equivalently the analytic implicit function theorem on this sequence
space, gives analytic dependence on $w_0$.  Evaluation at time zero
gives $\psi$.  The tangent space of a graph over $W_0$ has trivial
intersection with the pure thermal line, proving transversality.
\end{proof}

\begin{proposition}[The temperature branch]\label{flow:branch}
There is $\lambda_0>0$ and a real analytic branch
$\beta_c:(-\lambda_0,\lambda_0)\to(0,\infty)$, with
$\beta_c(0)=\beta_0$, whose homogeneous effective backgrounds satisfy
\begin{equation}\label{flow:plane-decay}
       h_{j+1}=\mathcal R_j(h_j),\qquad
       \norm{h_j}_{r_j}\le C|\lambda|\gamma_0^j.
\end{equation}
The branch depends only on $J,V,\lambda$.
\end{proposition}

\begin{proof}
Relative to the nearest-neighbor reference at $K_c$, the microscopic
background in the exponent is
\begin{equation}\label{flow:microscopic-background}
 (\beta J-K_c)\sum_{\{x,y\}:x\sim y}\sigma_x\sigma_y
       +\beta\lambda\sum_A V(A)\prod_{x\in A}\sigma_x.
\end{equation}
For the nearest-neighbor temperature displacement, assign half of each
unordered bond coefficient to each endpoint, and use its two-site carrier.
Thus its anchored function at $z$ is exactly
$2(\beta J-K_c)e(z)$. Assign finite connected carriers and anchors to the
remaining interaction shapes, distributing coefficients among equivalent
anchors when needed.  This preserves all lattice symmetries and spin-flip
parity.  These are lists of coefficients; no infinite-volume partition
function is being expanded.

Apply one preliminary version of \Cref{rg:map}, with protective
neighborhoods at radius $r_0$ and every microscopic singleton averaged
on its cutting square $B_{r_0}$.  The finite range makes the same
connected-list estimates convergent on a fixed neighborhood of
$(\beta_0,0)$; their constants may depend on $r_0,V$.  This gives an
analytic initial list $h_0(\beta,\lambda)\in\mathcal H_0$, equal to
zero at $(\beta_0,0)$.  The preliminary map is local, has the same
reserved exponential weight as the subsequent maps, and preserves
the finite-volume integral, with its scalar prefactor retained.

The normalization of its thermal derivative is explicit.  Because
each unoriented bond occurs twice in the sum of incident bonds,
\[
             \sum_z e(z)=\frac12\sum_{\{x,y\}:x\sim y}
                                           \sigma_x\sigma_y.
\]
Thus the preliminary derivative with respect to $\beta$, at
$(\beta_0,0)$, has the term $2J P_{r_0,z}e(z)$ at every anchor $z$.
By the definition of the homogeneous thermal list in
\Cref{rg:fields},
\begin{equation}\label{flow:thermal-derivative}
       \partial_\beta h_0(\beta_0,0)
                        =2Jr_0\boldsymbol v_0.
\end{equation}
This equality is in the quotient by constants and therefore also
holds for the canonical representatives.

Write $h_0(\beta,\lambda)=t_0(\beta,\lambda)\boldsymbol v_0
+w_0(\beta,\lambda)$.  The scalar analytic equation
\[
       t_0(\beta,\lambda)-\psi(w_0(\beta,\lambda))=0
\]
has $\beta$-derivative $2Jr_0\ne0$ at $(\beta_0,0)$, since
$\partial_\beta w_0(\beta_0,0)=0$.  The implicit function theorem
produces $\beta_c(\lambda)$.  All maps commute with complex
conjugation, so the branch and trajectory are real for real $\lambda$.
Shrinking the interval makes $\beta_c$ positive.  Along the branch,
$\norm{w_0}_{r_0}=O(|\lambda|)$, so \Cref{flow:stable-graph} gives
\eqref{flow:plane-decay}.  The argument applies to both signs of
$\lambda$ and uses no sign assumption on $V$.
\end{proof}

\subsection{The physical boundary}

Let $\Lambda_a=\{z\in\mathbb Z^2:az\in D\}$.  Call an ambient
carrier $A$ \emph{intact} if $A+B_1\subset\Lambda_a$, and let
$\operatorname{Res}_{a,j}h_j$ retain precisely the terms with intact
carriers.  All remaining finite-domain terms are represented by
connected ambient carriers meeting a fixed-width lattice neighborhood
of $\partial\Lambda_a$.  A carrier may extend outside $\Lambda_a$;
its function depends only on the available variable spins.  Terms
which are entirely scalar are removed from the background list.

\begin{proposition}[Boundary backgrounds]\label{flow:boundary}
There are constants $C,\lambda_0>0$, independent of the allowed domain
and of its boundary type, with the following property.  For each fixed
$C^2$ Jordan domain $D$ there is $\eta_D>0$ such that, for sufficiently
small $a$, the effective background with any of the three prescribed
boundary conditions has the form
\begin{equation}\label{flow:domain-decomposition}
       H_j^{a,b}=\operatorname{Res}_{a,j}h_j+b_j^{a,b},
       \qquad \norm{b_j^{a,b}}_{r_j}
                       \le C|\lambda|\gamma^j,
\end{equation}
through all steps with $ar_j\le\eta_D$.
Every carrier of $b_j^{a,b}$ touches the physical boundary in the
above sense.  For free boundary conditions its functions are even;
for fixed boundary conditions they need not be even.
\end{proposition}

\begin{proof}
A compact $C^2$ Jordan boundary has uniform one-sheet graph charts
whose slope variation is arbitrarily small on sufficiently short
physical windows.  Choose $\eta_D$ so that every enlarged local
window required by \Cref{rg:linear} is such a chart up to the stated
terminal scale.  The normalized chart constants are fixed universally;
only the permitted physical radius depends on $D$.  For small $a$
the preliminary scale also lies in these windows.

Before preliminary blocking, the finite-domain prescription differs
from the restriction of \eqref{flow:microscopic-background} only
within the fixed interaction range of the edge.  A crossing interaction
with fixed exterior spins remains a function of its available spins,
with an ambient carrier containing the original interaction set.
Deleting crossing interactions gives the corresponding free prescription.
Enlarging carriers by a fixed microscopic amount accommodates any
symmetry-preserving anchor decomposition.  The rooted norm of these
discrepancies is $O(|\lambda|)$, because
$\beta_c(\lambda)-\beta_0=O(|\lambda|)$.  The local preliminary map
therefore gives
\begin{equation}\label{flow:boundary-initial}
                    \norm{b_0^{a,b}}_{r_0}\le C_0|\lambda|.
\end{equation}
The bound uses local density of interaction shapes and the fixed
preliminary scale, not the total length of the boundary.

We verify the locality needed at subsequent steps.  For an intact
short-singleton output, its cutting square and all nearest-neighbor
edges required for conditional averaging are present in $\Lambda_a$.
The conditional kernel is exactly the plane kernel at that square.
For an intact connected output of order at least two, its carrier
contains all the protective neighborhoods of its inputs, and hence
every reference neighborhood used in its coefficient.  All these
inputs and kernels also agree with their plane counterparts.  A long
unchanged singleton uses no outside kernel.  Thus coefficients agree
on intact output carriers.  Any omission or discrepancy in comparing
the two maps has a carrier touching the edge; the same is true of
an output involving a boundary-family input.

The canonical representative subtracts the independent uniform-spin
mean on the available variables.  On an intact carrier that set of
variables is identical in the plane and the domain, so this subtraction
also agrees exactly.  At the edge the discarded scalar keeps its
generation carrier for any source bookkeeping.  Subtracting constants
therefore neither produces an interior discrepancy nor transports an
edge discrepancy into an interior carrier.

Let $\mathcal R_j^{a,b}$ denote the finite-domain map.  The preceding
locality and the absolute coefficient estimates of \Cref{rg:map}
give the bulk leakage bound
\begin{equation}\label{flow:boundary-leakage}
 \left\|\mathcal R_j^{a,b}(\operatorname{Res}_{a,j}h_j)
       -\operatorname{Res}_{a,j+1}\mathcal R_j(h_j)\right\|_{r_{j+1}}
                           \le C_L\norm{h_j}_{r_j}.
\end{equation}
No boundary contraction is asserted for this difference.  In contrast,
the linear image of the boundary family itself has norm at most
$\rho\norm{b_j^{a,b}}_{r_j}$ by \Cref{rg:linear}.  The nonlinear
terms involving that family cost at most
$C_L(\norm{h_j}_{r_j}+\norm{b_j^{a,b}}_{r_j})
\norm{b_j^{a,b}}_{r_j}$.  Consequently, writing
$q_j=\norm{b_j^{a,b}}_{r_j}$, we have
\begin{equation}\label{flow:boundary-recurrence}
 q_{j+1}\le\rho q_j+C_L\norm{h_j}_{r_j}
                     +C_L(\norm{h_j}_{r_j}+q_j)q_j.
\end{equation}
Evenness in the free case is preserved by every operation.  In the
fixed cases the unrestricted boundary estimate in \Cref{rg:linear}
applies instead.

Choose $\alpha$ with $\rho<\alpha<\gamma$.  On a sufficiently small
common ball the last term in \eqref{flow:boundary-recurrence} is at
most $(\alpha-\rho)q_j$.  Choose $B\ge C_0$ large enough that
$C_LC\le B(\gamma-\alpha)$, where $C$ is the constant in
\eqref{flow:plane-decay}.  Since $\gamma_0^j\le\gamma^j$, induction
then gives $q_j\le B|\lambda|\gamma^j$.  Decrease $\lambda_0$ so
this bound and \eqref{flow:plane-decay} stay in the analytic balls
and in the preceding smallness condition.  Every constant in this
choice is independent of $D$; the domain affects only $\eta_D$ and
the threshold for the mesh.  This closes the induction and proves
\eqref{flow:domain-decomposition}.
\end{proof}

\subsection{Periodic projection and exponentially small discrepancies}

Let $\mathbb T_N=(\mathbb Z/N\mathbb Z)^2$.  Fix a small numerical
$\kappa>0$, for example $\kappa=1/100$.  For a homogeneous plane
list, define $\Pi_N$ by retaining the carriers of diameter less than
$\kappa N$ and projecting their sites to $\mathbb T_N$.  Use one
representative of each anchor modulo $N\mathbb Z^2$, retaining its
relative carrier shapes and their multiplicities.  This projects a
homogeneous list once; it does not sum the already homogeneous list
over all its translates again.  Connected retained carriers lift
uniquely after their anchor lift is specified.  For $r\le\vartheta N$
with $\vartheta$ sufficiently small, comparison in a bounded number
of isometric charts gives
\begin{equation}\label{flow:projection-norm}
                       \norm{\Pi_N h}_r\le C\norm{h}_r.
\end{equation}

\begin{proposition}[Periodic backgrounds]\label{flow:torus}
For a sufficiently small fixed $\vartheta>0$, the periodic effective
background $H_j^N$ satisfies
\begin{align}
 \norm{H_j^N}_{r_j}&\le C|\lambda|\gamma^j,
                                      \label{flow:torus-small}\\
 \norm{H_j^N-\Pi_Nh_j}_{r_j}
       &\le C|\lambda|\gamma^j
                      \exp(-cN/r_j), \label{flow:torus-error}
\end{align}
whenever $r_j\le\vartheta N$.  The constants and the coupling
interval are independent of $N,j$.  On every fixed finite set of
scales, coefficients in a fixed plane chart converge to their plane
counterparts as $N\to\infty$, with the tails interpreted in the
weighted list norms.

For real $\lambda$ these backgrounds are real.  At every finite
volume the successive maps, including the preliminary map, retain
the exact partition identity with their scalar prefactors, and the
residual tilt defines a positive Gibbs probability.
\end{proposition}

\begin{proof}
Choose $\vartheta$ sufficiently small, depending on the already fixed
map parameters, that every retained carrier with its norm neighborhood
and every short-term cutting window under consideration lies in an
isometric torus chart.  The periodic microscopic interactions are
their usual images; only $N$ larger than the interaction range and
the preliminary scale is relevant here.

Compare one periodic map applied to $\Pi_Nh_j$ with
$\Pi_N\mathcal R_j(h_j)$.  A connected list for which the entire
union of protective and test neighborhoods lies in a sufficiently
small max-ball lifts uniquely, up to an overall period.  Its input
coefficients and contacts agree with those in the plane.  The local
conditional kernels on all singleton cutting squares agree as well,
and canonical uniform-spin subtraction uses identical variables.
The resulting coefficients therefore agree whenever neither an input
nor the output is removed by the diameter truncation.  This argument
uses equality of conditional kernels in a chart, not equality of
plane and torus marginal laws.

Every exception has extent at least a fixed multiple of $N$: it
contains an omitted input or output carrier, or its connected union
of test neighborhoods leaves every such small chart, including by
winding around the torus.  In the plane or torus where that list is
being estimated, the reserved exponential weight in \Cref{rg:map}
therefore supplies $\exp(-cN/r_{j+1})$.  Rooting near a norm block
costs only the usual local covering constant.  A long unchanged
singleton is compared directly with its truncated counterpart and
has the same tail bound.  Summing all exceptional lists, using the
small analytic ball to sum higher powers of the background, gives
\begin{equation}\label{flow:torus-forcing}
 \left\|\mathcal R_j^N(\Pi_Nh_j)
                   -\Pi_N\mathcal R_j(h_j)\right\|_{r_{j+1}}
        \le C_L\norm{h_j}_{r_j}\exp(-cN/r_{j+1}).
\end{equation}
The retained small carriers and their neighborhoods have comparable
norm weights in the two spaces by the chart choice.  The preliminary
finite-range map satisfies the same comparison, giving
\begin{equation}\label{flow:torus-initial}
              \norm{H_0^N-\Pi_Nh_0}_{r_0}
                    \le C_0|\lambda|\exp(-cN/r_0).
\end{equation}

Here is the summation that preserves a spatial exponential after
propagation.  Set $E_j=\norm{H_j^N-\Pi_Nh_j}_{r_j}$ and enlarge the
uniform local Lipschitz constant to $K\ge1$.  As long as the compared
backgrounds remain in the analytic balls, \eqref{flow:torus-forcing}
and \eqref{flow:plane-decay} imply
\[
 E_{j+1}\le K E_j+C_1|\lambda|\gamma^j
                                      \exp(-cN/r_{j+1}).
\]
For a terminal index $j$, put $x=N/r_j\ge\vartheta^{-1}$ and
$q=K/\gamma$.  Iteration, including \eqref{flow:torus-initial}, gives
\begin{equation}\label{flow:torus-convolution}
 \frac{E_j}{|\lambda|\gamma^j}
  \le C_0q^j e^{-cxL^j}
       +\frac{C_1}{\gamma}\sum_{m=0}^{j-1}q^m e^{-cxL^m}.
\end{equation}
At $\lambda=0$ the discrepancy is identically zero, so the displayed
division is needed only when $\lambda\ne0$.  For $0<c'<c$,
\[
 \sum_{m\ge0}q^m e^{-cxL^m}
 \le e^{-c'x}
    \sum_{m\ge0}q^m
           \exp\bigl(-(cL^m-c')/\vartheta\bigr)
 \le C_{L,K,\gamma,\vartheta}e^{-c'x}.
\]
The series converges because its exponential decay in $L^m$
dominates $q^m$.  The initial term in
\eqref{flow:torus-convolution} is its $m=j$ term and has the same
bound.  This proves \eqref{flow:torus-error}, after renaming $c'$.
Together with \eqref{flow:projection-norm} it gives
\eqref{flow:torus-small}.  Decreasing the common $\lambda_0$ makes
these bounds stay strictly inside all analytic balls, and closes the
finite-scale induction uniformly in $N$.

For fixed $j$, \eqref{flow:torus-error} tends to zero as
$N\to\infty$.  The omitted plane carriers also have vanishing
weighted tails at that fixed scale.  Identifying the remaining local
carriers with their unique lifts proves the claimed coefficient
convergence.  Finally the preliminary and subsequent maps preserve
real coefficients and the exact finite-volume identity of
\Cref{rg:map}.  Exponentiating the real residual background gives
a strictly positive weight, completing the last assertion.
\end{proof}
